\documentclass[english, parskip=full]{scrartcl}
\usepackage[english]{babel}  % Sprache auf Deutsch 
\usepackage[utf8]{inputenc}  % Kodierung der Datei
\usepackage[left=2cm, right=2cm, top=2cm, bottom=3cm, footskip=1.5cm]{geometry}
\usepackage{color}
\usepackage{graphicx, gensymb}
\usepackage{amsfonts, cancel, amssymb}
\usepackage{amsmath, amsthm, array, esint}
\usepackage{bm} % bold math symbols, e.g. \bm{\sigma} etc.
\usepackage{booktabs, multirow}  % schönere Tabellen
%\usepackage{scrlayer-scrpage}    % Manipulation des Seitenstils
%\pagestyle{scrheadings}  % Stil für die Seite setzen
\usepackage{tabularx}
\usepackage{setspace}
%\automark{section}  % Abschnittsnamen als Seitenbeschriftung 
%\setheadsepline{.5pt}    % Kopzeile durch Linie abtrennen
\usepackage[hidelinks]{hyperref}  % Links und weitere PDF-Features
\usepackage{typearea, tikz, pgffor, verbatim}
\usetikzlibrary{calc, arrows.meta,arrows, graphs, quotes, positioning, angles, babel}
\usepackage[cache=false, outputdir=build]{minted}
\usemintedstyle{vs}
\usepackage{placeins} % provides \FloatBarrier

\definecolor{bg}{rgb}{0.9,0.9,0.9}
\newcommand\numberthis{\addtocounter{equation}{1}\tag{\theequation}}
\usepackage{svg}
\usepackage{siunitx}
\usepackage{physics}
\usepackage{tensor}
\usepackage{slashed}
\usepackage{bbm}
\usepackage[compat=1.1.0]{tikz-feynman}

\usepackage{caption}
\captionsetup{
    % width=.8\textwidth,
    % indention=1cm,
    margin=0.5cm,
    format=plain,
    labelfont=bf
}
%\usepackage{subcaption}
%\subcaptionsetup{
%    indention=0cm,
%    format=hang,
%    margin=0.5cm,
%    labelfont=normalfont
%}
\usepackage[english,capitalise]{cleveref}
\usepackage[
    backend=biber,
    sorting=none,
    style=phys,
    giveninits=true,
    sortcites=true,
    citestyle=numeric-comp,
    % url=false,
    % doi=false,
    biblabel=brackets,
    % eprint=false,
    maxcitenames=3]{biblatex}
\usepackage{csquotes}
%\addbibresource{bibliography.bib}

\usepackage{mathtools}
% use \abs for absolute values
% use \abs for \left ... \right version and \abs[\bigg for \bigg version etc.
\let\abs\undefined
\let\bra\undefined
\let\braket\undefined
\DeclarePairedDelimiter\abs{\vert}{\vert}
\DeclarePairedDelimiter\kla{(}{)}
\DeclarePairedDelimiter\klb{[}{]}
\DeclarePairedDelimiter\klc{\{}{\}}
\DeclarePairedDelimiter\bra{\langle}{\rangle}
\DeclarePairedDelimiterX\brb[2]{\langle #1}{\rangle}{\delimsize\vert #2 }
\DeclarePairedDelimiterX\brc[3]{\langle #1}{\rangle}{\delimsize\vert #2 \delimsize\vert #3 }

\renewcommand{\i}{\text{i}}
\newcommand{\e}{\text{e}}
\DeclareMathOperator{\sgn}{sgn}
\newcommand{\kom}[1]{\overset{\substack{#1 \\ |}}{=}}
\newcommand{\koml}[2]{\overset{\substack{#1 \\ |}}{#2}}
\newcommand{\intx}[4]{\int_{#1}^{#2} #3 \,\mathrm{d}#4}
\newcommand{\intr}[2]{\int_{\mathbb{R}} #1 \, \mathrm{d}#2}
\newcommand{\intrnd}[3]{\int_{\mathbb{R}^{#1}} #2 \, \mathrm{d}^{#1}#3}
\newcommand{\intrpi}[2]{\int_{\mathbb{R}} #1 \, \frac{\mathrm{d}#2}{2\pi}}
\newcommand{\intrndpi}[3]{\int_{\mathbb{R}^{#1}} #2 \, \frac{\mathrm{d}^{#1}#3}{(2\pi)^{#1}}}
\newcommand{\oper}[2]{\vert #1 \rangle \langle #2 \vert}
\newcommand{\Text}[1]{\text{#1}}    % this is relevant because \text does not autocomplete to the default '\text{@}' but rather '\textbf' etc.
\newcommand{\powe}[1]{\e^{#1}}
\newcommand{\pow}[2]{{#1}^{#2}}
\newcommand{\Ket}[1]{\vert #1 \rangle}
\newcommand{\Bra}[1]{\langle #1 \vert}
\newcommand*{\Fourier}{\textsc{Fourier}\ }
\newcommand*{\F}{\mathcal{F}}
\DeclareMathOperator{\arsinh}{arsinh}
\DeclareMathOperator{\arcosh}{arcosh}
\DeclareMathOperator{\artanh}{artanh}
\DeclareMathOperator{\sinc}{sinc}
\DeclareMathOperator{\cscc}{cscc}
\DeclareMathOperator{\sinhc}{sinhc}
\DeclareMathOperator{\cschc}{cschc}
\newcommand{\conv}{\mathop{\scalebox{3.5}{\raisebox{-0.38ex}{\makebox[0.5\width][c]{$\ast$}}}}\limits}
\newcommand*{\unity}{\text{\usefont{U}{bbold}{m}{n}1}}   % operator one from :: https://tex.stackexchange.com/questions/566780/import-mathbb1-character-from-the-unicode-math-package

\renewcommand{\d}{\text{d}}
\renewcommand{\r}{\text{r}}
\renewcommand{\t}{\text{t}}
\renewcommand{\a}{\text{a}}
\renewcommand{\o}{\text{o}}
\renewcommand{\F}{\text{F}}
\renewcommand{\c}{\text{c}}
\newcommand{\A}{\text{A}}
\newcommand{\B}{\text{B}}
\newcommand{\R}{\text{R}}
\newcommand{\M}{\text{M}}
\newcommand{\m}{\text{m}}
\newcommand{\s}{\text{s}}
\newcommand{\f}{\text{f}}
\newcommand{\K}{\text{K}}
\newcommand{\hc}{\text{h.c.}}
\newcommand{\kf}{k_\F}
\newcommand{\vf}{v_\F}
\newcommand{\const}{\text{const.}}
\renewcommand{\L}{\text{L}}
\newcommand{\gray}[1]{\bgroup\color{white!40!black}#1\egroup}
\newcommand{\TODO}[1]{\bgroup\color{red!60!black}#1\egroup}
\newcommand\QnA[1]{\bgroup\color{blue}#1\egroup}
\newcommand\highlight[1]{\bgroup\color{green!60!black}#1\egroup}

\newcommand{\cabx}[3]{\begin{smallmatrix}\gray{A}&\gray{:}&#1 \\ \gray{B}&\gray{:}&#2 \\ \gray{\Xi}&\gray{:}&#3\end{smallmatrix}}
\newcommand{\zabx}[3]{\begin{smallmatrix}\gray{A}&\gray{:}&#1 \\ \gray{B}&\gray{:}&#2 \\ \gray{Z}&\gray{:}&#3\end{smallmatrix}}
\newcommand{\abx}[3]{\begin{smallmatrix}#1 \\ #2 \\ #3\end{smallmatrix}}

\newcommand{\normord}[1]{
  {:\mathrel{\mspace{1mu}#1\mspace{1mu}}:}
}
\newcommand{\varnormord}[1]{
  {\mathrel{\mspace{1mu}\substack{\ast\\[-1.5pt] \ast}#1\substack{\ast\\[-1.5pt] \ast}\mspace{1mu}}}
}

% punctuation styles (see https://tex.stackexchange.com/questions/56063/how-to-properly-typeset-footnotes-superscripts-after-punctuation-marks)
\let\origfootnote\footnote
\renewcommand{\footnote}[1]{\kern.06em\origfootnote{#1}}
\newcommand{\punctfootnote}[1]{\kern-.06em\origfootnote{#1}}

% Multiline equations
\newcommand{\bsplit}[1]{\begin{aligned}[t]#1\end{aligned}}

\newcommand*{\titel}{Green's functions for simple systems}
\newcommand*{\autor}{Joachim Schwardt}

\hypersetup{pdfauthor={\autor}, pdftitle={\titel}}

%\titlehead{titlehead}
\subject{Condensed Matter Physics}
\title{\titel}
\author{\autor}
\date{\today}

\begin{document}

%\begin{titlepage}
\maketitle  % Titel setzen
%\tableofcontents  % Inhaltsverzeichnis setzen
%\end{titlepage}

\section{Free fermionic systems}
\subsection{Fermionic language}
Consider the non-interacting Hamiltonian in momentum space
\begin{align}
H &= \sum_{k\sigma} \epsilon_{k\sigma} c_{k\sigma}^\dagger c_{k\sigma}
\end{align}
with the dispersion relation $\epsilon_{k\sigma}$ and the fermionic creation and annihilation operators $c_{k\sigma}^\dagger$ and $c_{k\sigma}$.
Due to the Heisenberg equation the operators pick up a time dependence according to\footnote{$\klb*{a, bc} = \klc*{a,b}c - b\klc*{a,c}$}
\begin{align}
\i\partial_t c_{k\sigma} &= \klb*{c_{k\sigma}, H} = \sum_{k'\sigma'} \epsilon_{k'\sigma'} \klb*{c_{k\sigma}, c_{k'\sigma'}^\dagger c_{k'\sigma'}} = \epsilon_{k\sigma} c_{k\sigma}.
\end{align}
This equation yields $c_{k\sigma}(t) = c_{k\sigma} \powe{-\i\epsilon_{k\sigma} t}$ and hermitian conjugation gives $c_{k\sigma}^\dagger(t) = c_{k\sigma}^\dagger \powe{\i\epsilon_{k\sigma} t}$.
The retarded and advanced Green's functions are defined as
\begin{align}
G^\r_{\sigma\sigma'}(x,t;x',t') &= -\i\Theta(t-t') \bra*{\klc*{c_{x\sigma}(t), c_{x'\sigma'}^\dagger(t')}},\\
G^\a_{\sigma\sigma'}(x,t;x',t') &= \i\Theta(t'-t) \bra*{\klc*{c_{x\sigma}(t), c_{x'\sigma'}^\dagger(t')}}.
\end{align}
Performing a spatial Fourier transform in both arguments gives
\begin{align}
G^\r_{\sigma\sigma'}(k,t;k',t') &= \intr{\intr{\powe{-\i kx} \powe{\i k'x'} G^\r_{\sigma\sigma'}(x,t;x',t') }{x'}}{x} \\
&= -\i\Theta(t-t') \bra*{\klc*{c_{k\sigma}, c_{k'\sigma'}^\dagger}} \powe{-\i\epsilon_{k\sigma}t} \powe{\i\epsilon_{k'\sigma'}t'} = -\i\Theta(t-t') \delta_{kk'} \delta_{\sigma\sigma'} \powe{-\i\epsilon_{k\sigma} (t-t')},\\
G^\a_{\sigma\sigma'}(k,t;k',t') &= \i\Theta(t'-t) \delta_{kk'} \delta_{\sigma\sigma'} \powe{-\i\epsilon_{k\sigma} (t-t')}.
\end{align}
The diagonal structure allows us to lighten the notation a bit by writing
\begin{align}
G^\r_\sigma(k,t) &= -\i\Theta(t) \powe{-\i\epsilon_{k\sigma} t},\\
G^\a_\sigma(k,t) &= \i\Theta(-t) \powe{-\i\epsilon_{k\sigma} t}.
\end{align}
The temporal Fourier transform now yields
\begin{align}
G^\r_\sigma(k,\omega) &= \intr{\powe{\i\omega t} G^\r_\sigma(k,t)}{t} = -\i\lim_{\eta\rightarrow 0} \intx{0}{\infty}{\powe{\i\omega t}\powe{-\i\epsilon_{k\sigma} t} \powe{-\eta t}}{t} = \frac{1}{\omega - \epsilon_{k\sigma} + \i 0^+},\\
G^\a_\sigma(k,\omega) &= \intr{\powe{\i\omega t} G^\r_\sigma(k,t)}{t} = \i\lim_{\eta\rightarrow 0} \intx{-\infty}{0}{\powe{\i\omega t}\powe{-\i\epsilon_{k\sigma} t} \powe{\eta t}}{t} = \frac{1}{\omega - \epsilon_{k\sigma} - \i 0^+}.
\end{align}
These results are connected to the Matsubara Green's function which is defined by
\begin{align}
G^\M_{\sigma\sigma'}(x,\tau;x',\tau') &= -\bra*{\mathcal{T}_\tau c_{x\sigma}(\tau) c_{x'\sigma'}^\dagger(\tau')}.
\end{align}
As before the spatial transform gives
\begin{align}
G^\M_{\sigma\sigma'}(k,\tau;k',\tau') &= -\bra*{\mathcal{T}_\tau c_{k\sigma}(\tau) c_{k'\sigma'}^\dagger(\tau')} = -\bra*{c_{k\sigma} c_{k'\sigma'}^\dagger} \powe{-\epsilon_{k\sigma} \tau} \powe{\epsilon_{k'\sigma'} \tau'}.
\end{align}
where we assumed the imaginary time $\tau = \i t$ to be larger than $\tau'$.
The expectation values now refer to thermal averages
\begin{align}
\bra*{c_{k\sigma} c_{k'\sigma'}^\dagger} &= \frac{1}{Z} \tr\klc*{\powe{-\beta H} c_{k\sigma} c_{k'\sigma'}^\dagger} = \delta_{kk'} \delta_{\sigma\sigma'} - \frac{1}{Z}\tr\klc*{\powe{-\beta H} c_{k'\sigma'}^\dagger c_{k\sigma}}.
\end{align}
Using the cyclic property of the trace and\footnote{$\klb*{n_{k\sigma}, n_{k'\sigma'}} = 0$ and $\klb*{n_{k\sigma}, c_{k'\sigma'}} = -\delta_{kk'}\delta_{\sigma\sigma'} c_{k\sigma}$}
\begin{align}
c_{k\sigma} \powe{-\beta H} &= c_{k\sigma} \prod_{k'\sigma'} \powe{-\beta \epsilon_{k'\sigma'} n_{k'\sigma'}} = \powe{-\beta H} \powe{\beta \epsilon_{k\sigma} n_{k\sigma}} c_{k\sigma} \powe{-\beta \epsilon_{k\sigma} n_{k\sigma}} \\
&= \powe{-\beta H} \sum_{n\ge 0} \frac{1}{n!} \klb*{\beta \epsilon_{k\sigma} n_{k\sigma}, c_{k\sigma}}_n = \powe{-\beta H} \sum_{n\ge 0} \frac{(-1)^n}{n!} (\beta\epsilon_{k\sigma})^n c_{k\sigma} = \powe{-\beta H} \powe{-\beta \epsilon_{k\sigma}} c_{k\sigma}
\end{align}
we find the relation
\begin{align}
\bra*{c_{k\sigma} c_{k'\sigma'}^\dagger} &= \delta_{kk'} \delta_{\sigma\sigma'} - \frac{1}{Z} \tr\klc*{\powe{-\beta H} \powe{-\beta\epsilon_{k\sigma}} c_{k\sigma} c_{k'\sigma'}^\dagger} = \delta_{kk'} \delta_{\sigma\sigma'} - \powe{-\beta\epsilon_{k\sigma}} \bra*{c_{k\sigma} c_{k'\sigma'}^\dagger}.
\end{align}
Rearranging yields the well-known Fermi-Dirac statistics
\begin{align}
\bra*{c_{k\sigma} c_{k'\sigma'}^\dagger} &= \frac{\delta_{kk'} \delta_{\sigma\sigma'}}{\powe{-\beta\epsilon_{k\sigma}} + 1} = \delta_{kk'} \delta_{\sigma\sigma'} \kla*{1 - \frac{1}{\powe{\beta \epsilon_{k\sigma}} + 1}} = \delta_{kk'} \delta_{\sigma\sigma'} \kla*{1 - n^\F_{k\sigma}}.
\end{align}
Once again the Green's function is diagonal in $k$ and $\sigma$ so we abbreviate by writing
\begin{align}
G^\M_\sigma(k,\tau) &= -\kla*{1 - n_{k\sigma}^\F} \powe{-\epsilon_{k\sigma} \tau}.
\end{align}
To get the frequency representation we compute the Fourier transform with the fermionic Matsubara frequencies $\omega_n^\F = \frac{(2n+1)\pi}{\beta}$
\begin{align}
G^\M_\sigma(k,\i\omega_n^\F) &= \intx{0}{\beta}{\powe{\i\omega_n^\F \tau} G^\M_\sigma(k,\tau)}{\tau} = \kla*{n^\F_{k\sigma} - 1}\intx{0}{\beta}{\powe{\i\omega_n^\F \tau - \epsilon_{k\sigma} \tau}}{\tau} \\
&= \kla*{n^\F_{k\sigma} - 1} \frac{\powe{\i\omega_n^\F \beta} \powe{-\beta \epsilon_{k\sigma}} - 1}{\i\omega_n^\F - \epsilon_{k\sigma}} = \frac{1}{\i\omega_n^\F - \epsilon_{k\sigma}}.
\end{align}
Thus we see that the retarded and advanced Green's function can be obtained from the Matsubara version by performing the analytic continuation $\i\omega_n^\F \rightarrow \omega \pm \i 0^+$.
\subsection{Bosonized language}
We now want to demonstrate that all of the above calculations can also be done within the formalism of Bosonization assuming that we are dealing with a linear dispersion $\epsilon_{k\sigma} = \vf\sigma k$, where $\sigma = \pm 1$ are the two possible states.
In real space we make use of the operator identity
\begin{align}
c_{x\ell} &= \sqrt{a} \psi_\ell(x) = \sqrt{a} \frac{\eta_\ell}{\sqrt{2\pi\alpha}} \powe{-\i\varphi_\ell(x)}
\end{align}
which gives rise to the quadratic Hamiltonian
\begin{align}
H &= \frac{\vf}{2\pi} \intr{\klb*{\kla*{\partial_x\theta}^2 + \kla*{\partial_x\phi}^2}}{x}
\end{align}
where $\vf$ is the Fermi velocity -- which we will now set to unity -- and $\varphi_\ell(x) = \ell \phi(x) - \theta(x)$.
Note that the fields $\phi$ and $\theta$ commute among themselves and satisfy
\begin{align}
\klb*{\phi(x),\theta(y)} &= \i\frac{\pi}{2} \sgn(x-y) = \klb*{\theta(x), \phi(y)}.
\end{align}
As before we start with the time evolution of the operators in momentum space
\begin{align}
\i\partial_t \psi_\ell(k,t) &= \intr{\powe{-\i kx} \i\partial_t\psi_\ell(x,t)}{x} = \intr{\powe{-\i kx} \klb*{\psi_\ell(x,t), H}}{x}\\
&= \intr{\powe{-\i kx} \frac{1}{2\pi}\intr{ \klb*{\kla*{\partial_y\theta}^2 + \kla*{\partial_y\phi}^2, \psi_\ell(x,t)} }{y} }{x}\\
&= \frac{1}{\pi} \intr{ \powe{-\i kx} \intr{ \klb[\big]{ \partial_y\klb*{\theta, \psi_\ell} \partial_y\theta + \partial_y\phi \klb*{\phi, \psi_\ell} } }{y} }{x}\\
&= \frac{1}{\pi} \intr{\powe{-\i kx} \intr{ \klb*{\i\pi\delta(y-x)(-\i\ell) \psi_\ell \partial_y\theta + \i\pi\delta(y-x) \i \psi_\ell \partial_y\phi} }{y} }{x}\\
&= -\ell \intr{\powe{-\i kx} \psi_\ell \partial_x\varphi_\ell}{x} = -\ell\i\intr{\powe{-\i kx} \partial_x\psi_\ell}{x} = \ell k \psi_\ell(k,t).
\end{align}
Solving this differential equation yields $\psi_\ell(k,t) = \psi_\ell(k) \powe{-\i\ell kt} \equiv \psi_\ell(k) \powe{-\i \epsilon_{\ell k} t}$.\\
To proceed with the calculation of the Green's function we now need to consider the anti-commutator $\klc*{c_{k\ell}, c_{k',\ell'}^\dagger} = \delta_{kk'} \delta_{\ell\ell'}$.
In our bosonized language this corresponds to\footnote{Note that the notation is a bit lose here because we do not pay close attention to discrete or continuous momenta. 
The qualitative description is correct though, and the details are identical in the fermionic and the bosonic language.}
\begin{align}
\klc*{\psi_\ell(k), \psi_{\ell'}^\dagger(k')} &= \intr{\intr{\powe{-\i kx} \powe{\i k'x'} \klc*{\psi_\ell(x), \psi_{\ell'}^\dagger(x')} }{x'} }{x}\\
&= \delta_{\ell\ell'} \intr{\powe{-\i x(k-k')}}{x} = \delta_{\ell\ell'} \delta_{kk'}.
\end{align}
Here we made use of $\klc*{\eta_\ell, \eta_{\ell'}} = \delta_{\ell\ell'}$ and of $\klb*{\varphi_\ell^{(+)}(x), \varphi_{\ell'}^{(-)}(x')} = \delta_{\ell\ell'} \log\klb*{\frac{2\pi}{L} \kla*{\alpha + \i\ell (x-x')}}$ to prove the anti-commutation relation\footnote{$\psi_\ell(x) = \frac{\eta_\ell}{\sqrt{2\pi\alpha}} \powe{-\i\varphi_\ell^{(+)} - \i\varphi_\ell^{(-)}} = \frac{\eta_\ell}{\sqrt{2\pi\alpha}} \normord{\powe{-\i\varphi_\ell}} \powe{-\frac{1}{2}\klb*{-\i\varphi_\ell^{(+)}, -\i\varphi_\ell^{(-)}} } = \frac{\eta_\ell}{\sqrt{L}} \normord{\powe{-\i\varphi_\ell}}$}
\begin{align}
\klc*{\psi_\ell(x), \psi_{\ell'}^\dagger(x')} &= \delta_{\ell\ell'} \frac{1}{L} \klc*{\normord{\powe{-\i\varphi_\ell(x)}}, \normord{\powe{\i \varphi_\ell(x')}}} \\
&= \delta_{\ell\ell'} \frac{1}{L} \normord{\powe{-\i\varphi_\ell(x) + \i\varphi_\ell(x')}} \kla*{\powe{\klb*{-\i\varphi_\ell^{(-)}(x), \i\varphi_\ell^{(+)}(x')}} + \powe{\klb*{\i\varphi_\ell^{(+)}(x'), -\i\varphi_\ell^{(-)}(x)}}}\\
&= \delta_{\ell\ell'} \frac{1}{L} \normord{\powe{-\i\varphi_\ell(x) + \i\varphi_\ell(x')}} \kla*{\frac{L}{2\pi} \frac{1}{\alpha - \i\ell(x-x')} + \frac{L}{2\pi} \frac{1}{\alpha + \i\ell(x-x')}}\\
&= \delta_{\ell\ell'} \normord{\powe{-\i\varphi_\ell(x) + \i\varphi_\ell(x')}} \frac{1}{\pi} \frac{\alpha}{\alpha^2 + (x-x')^2} = \delta_{\ell\ell'} \delta(x-x').
\end{align}
\section{Interactions and perturbation theory}
In this section we want to consider the same Hamiltonian as before, now denoted as $H_0$, and a simple bilinear interaction term $H_\Text{int} = \sum_{k\sigma} g c_{k\sigma}^\dagger c_{k,-\sigma}$.
To keep calculations simple we assume a linear dispersion $\epsilon_{k\sigma} = \sigma k$.
\subsection{Fermionic language}
\subsubsection{Exact solution}
Obviously this problem is exactly soluble due to the quadratic nature of the Hamiltonian.
All we need to do is to diagonalize the matrix $M = k\sigma_z + g\sigma_x$ and one can check that
\begin{align}
M &= UDU^T
\end{align}
with $D = E\sigma_z$, $E=\sqrt{k^2 + g^2}$ and the transformation matrix
\begin{align}
U &= \frac{1}{\sqrt{2E(E-k)}} \begin{pmatrix}
g & k-E \\ E-k & g
\end{pmatrix}
\end{align}
%from green_toolkit import sigma_0, sigma_x, sigma_y, sigma_z
%k=0.43; g=0.114
%H=k*sigma_z + g*sigma_x
%E=np.sqrt(k**2+g**2)
%U=(g*sigma_0 + 1j*sigma_y*(k-E))/np.sqrt(2*E*(E-k))
%D=E*sigma_z
%H, U@D@U.T
is the correct decomposition.
This also gives rise to new operators
\begin{align}
d_{k\sigma} &= \frac{1}{\sqrt{2E(E-k)}} \kla*{gc_{k\sigma} + \sigma (E-k) c_{k,-\sigma}},\\
c_{k\sigma} &= \frac{1}{\sqrt{2E(E-k)}} \kla*{gd_{k\sigma} - \sigma (E-k) d_{k,-\sigma}}
\end{align}
The full hamiltonian can then be written as $H=H_0+H_\Text{int} = \sum_{k\sigma} \sigma E d_{k\sigma}^\dagger d_{k\sigma}$.
The retarded Green's function becomes
\begin{align}
G_{\sigma\sigma'}^\r(k,t) &= -\i\Theta(t) \bra*{\klc*{c_{k\sigma}(t) , c_{k\sigma'}^\dagger(0)}} \\
&= -\i\Theta(t) \frac{1}{2E(E-k)} \bra*{ \klc*{gd_{k\sigma} \powe{-\i \sigma Et} - \sigma (E-k)d_{k,-\sigma} \powe{\i \sigma Et} , gd_{k\sigma'}^\dagger -\sigma' (E-k) d_{k,-\sigma'}^\dagger} }\\
&= -\i\Theta(t) \frac{1}{2E(E-k)} \klc*{\delta_{\sigma\sigma'} \klb*{g^2 \powe{-\i\sigma Et} + (E-k)^2 \powe{\i\sigma Et}} - \delta_{\sigma,-\sigma'} g (E-k) 2\i\sin(Et)}.
\end{align}
The Fourier transform yields
\begin{align}
G_{\sigma\sigma}^\r(k,\omega) &= -\i \intx{0}{\infty}{\powe{\i\omega t} \powe{-\delta t} G_{\sigma\sigma'}^\r(k,t) }{t}\\
&= \frac{1}{2E(E-k)} \klc*{\delta_{\sigma\sigma'} \klb*{\frac{g^2}{\omega - \sigma E + \i\delta} + \frac{(E-k)^2}{\omega + \sigma E + \i\delta}} - \delta_{\sigma,-\sigma'} \frac{2E g(E-k)}{E^2 - \omega^2 - 2\i\omega \delta}}\\
&= \frac{1}{\omega^2 - E^2} \klc*{\delta_{\sigma\sigma'} \klb*{\omega + \sigma k} + \delta_{\sigma,-\sigma'} g}.
\end{align}
We can obtain the same result within the Matsubara formalism
\begin{align}
G_{\sigma\sigma'}^\M(k,t) &= -\mathcal{T} \bra*{c_{k\sigma}(\tau) c_{k\sigma'}^\dagger(0)} \\
&= -\frac{1}{2E(E-k)} \bra*{ \kla*{gd_{k\sigma} \powe{- \sigma E\tau} - \sigma (E-k)d_{k,-\sigma} \powe{\sigma E\tau} } \kla*{gd_{k\sigma'}^\dagger -\sigma' (E-k) d_{k,-\sigma'}^\dagger} }\\
&= -\frac{1}{2E(E-k)} \bsplit{\klc[\bigg]{&\delta_{\sigma\sigma'} \klb*{g^2 \powe{-\sigma E\tau} \frac{1}{1 + \powe{-\sigma E\beta}} + (E-k)^2\powe{\sigma E\tau} \frac{1}{1 + \powe{\sigma E\beta}} } \\
&+\delta_{\sigma,-\sigma'} \sigma g(E-k) \klb*{\powe{-\sigma E\tau} \frac{1}{1 + \powe{-\sigma E\beta}} - \powe{\sigma E\tau} \frac{1}{1 + \powe{\sigma E\beta}}} }.}
\end{align}
The Fourier transform for the diagonal part yields
\begin{align}
G_{\sigma\sigma}^\M(k,\i\omega_n^\F) &= \intx{0}{\beta}{\powe{\i\omega_n^\F\tau} G_{\sigma\sigma}^\M(k,\tau)}{\tau}\\
&= \frac{-1}{2E(E-k)} \klc*{\frac{g^2}{\i\omega_n^\F - \sigma E} \frac{-\powe{-\sigma E\beta} - 1}{1 + \powe{-\sigma E\beta}} + \frac{(E-k)^2}{\i\omega_n^\F + \sigma E} \frac{-\powe{\sigma E\beta} - 1}{1 + \powe{\sigma E\beta}}}\\
&= \frac{1}{\kla*{\i\omega_n^\F}^2 - E^2} \klc*{\i\omega_n^\F + \sigma k}
\end{align}
and the one for the off-diagonal part gives
\begin{align}
G_{\sigma,-\sigma}^\M(k,\i\omega_n^\F) &= \intx{0}{\beta}{\powe{\i\omega_n^\F\tau} G_{\sigma,-\sigma}^\M(k,\tau)}{\tau}\\
&= \frac{-\sigma g}{2E} \klc*{\frac{-1}{\i\omega_n^\F - \sigma E} - \frac{-1}{\i\omega_n^\F + \sigma E}} = \frac{g}{\kla*{\i\omega_n^\F}^2 - E^2}.
\end{align}
Performing the analytic continuation $\i\omega_n^\F \rightarrow \omega + \i\delta$ reproduces the result from the real-time calculation.
Note that the Matsubara approach was quite a bit more cumbersome.
\subsubsection{Perturbation theory}
Taylor expanding the exact result to first order in $g$ gives
\begin{align}
G_{\sigma\sigma'}^\r(k,\omega) &\simeq \frac{\omega + k\sigma_z}{\omega^2 - k^2} + \frac{g\sigma_x}{\omega^2 - k^2} \equiv G_{0,\sigma\sigma'}^\r(k,\omega) - g G_{1,\sigma\sigma'}^\r(k,\omega).
\end{align}
The additional minus sign originates from the Taylor expansion of $\powe{-S_\Text{int}}$, which is the weighting function appearing the the path integral in imaginary time.
The free Green's function can be obtained in a number of ways, some of which shown in the first section.
We thus only need to understand how the correction term comes about.
This time let us start with the Matsubara formalism and note that $\bra*{S_\Text{int}}_0 = 0$ due to the off-diagonal structure.
Therefore the first order correction should be
\begin{align}
G_{1,\sigma\sigma'}^\M(k,t) &= - \sum_{qs} \intx{0}{\beta} {\mathcal{T} \bra*{c_{k\sigma}(\tau) c_{k\sigma'}^\dagger(0) c_{qs}^\dagger(\tau') c_{q,-s}(\tau')}_0}{\tau'}\\
&= -\delta_{\sigma,-\sigma'} \powe{-\sigma k\tau} \klb[\bigg]{\bra*{c_{k\sigma} c_{k\sigma}^\dagger c_{k,-\sigma} c_{k,-\sigma}^\dagger}_0 \intx{0}{\tau}{ \powe{2\sigma k\tau'} }{\tau'} + \bra*{c_{k\sigma}^\dagger c_{k,-\sigma} c_{k\sigma} c_{k,-\sigma}^\dagger}_0 \intx{\tau}{\beta}{ \powe{2\sigma k\tau'} }{\tau'}}\\
&= -\delta_{\sigma,-\sigma'} \powe{-\sigma k\tau} \klb*{\bra*{\kla*{1 - \hat{n}_{k\sigma}} \kla*{1 - \hat{n}_{k,-\sigma}}}_0 \frac{\powe{2\sigma k\tau} - 1}{2\sigma k} - \bra*{\hat{n}_{k\sigma} \kla*{1 - \hat{n}_{k,-\sigma}}}_0 \frac{\powe{2\sigma k\beta} - \powe{2\sigma k\tau}}{2\sigma k}}.
\end{align}
The expectation values are simple to evaluate if one realizes that
\begin{align}
\bra*{\hat{n}_{k\sigma} \hat{n}_{k,-\sigma}}_0 &= \frac{1}{\mathcal{Z}} \sum_{\klc*{n_{qs}}} \brc*{\klc*{n_{qs}}}{\hat{n}_{k\sigma} \hat{n}_{k,-\sigma} \powe{-\beta H_0}}{\klc*{n_{qs}}}\\
&= \frac{1}{\mathcal{Z}_k} \sum_{n_{k,+}=0}^1 \sum_{n_{k,-}=0}^1 n_{k\sigma} n_{k,-\sigma} \powe{-\beta k n_{k,+} + \beta k n_{k,-}} = \bra*{\hat{n}_{k\sigma}}_0\bra*{\hat{n}_{k,-\sigma}}_0.
\end{align}
Using $\bra*{\hat{n}_{k\sigma}}_0 = \frac{1}{\powe{\sigma k\beta} + 1}$ and $\bra*{1 - \hat{n}_{k\sigma}} = \frac{1}{1 + \powe{-\sigma k\beta}}$ we then find
\begin{align}
G_{1,\sigma\sigma'}^\M(k,t) &=-\delta_{\sigma,-\sigma'} \powe{-\sigma k\tau} \klb*{\frac{1}{1 + \powe{-\sigma k\beta}} \frac{1}{1 + \powe{\sigma k\beta}} \frac{\powe{2\sigma k\tau} - 1}{2\sigma k} - \frac{1}{\powe{\sigma k\beta} + 1} \frac{1}{1 + \powe{\sigma k\beta}} \frac{\powe{2\sigma k\beta} - \powe{2\sigma k\tau}}{2\sigma k}}\\
&= -\delta_{\sigma,-\sigma'} \frac{1}{4 k\cosh^2\kla*{\frac{k\beta}{2}}} \klb*{\sinh(k\tau) - \sinh(k\beta - k\tau)}.
\end{align}
Performing the Fourier transform yields
\begin{align}
G_{1,\sigma\sigma'}^\M(k,\i\omega_n^\F) &= \frac{-\delta_{\sigma,-\sigma'}}{4k\cosh^2\kla*{\frac{k\beta}{2}}} \intx{0}{\beta}{ \klb*{ \powe{\i\omega_n^\F \tau} \sinh(k\tau) - \powe{\i\omega_n^\F\tau} \sinh(k\beta - k\tau)}}{\tau}\\
&= \frac{-\delta_{\sigma,-\sigma'}}{4k\cosh^2\kla*{\frac{k\beta}{2}}} \frac{1}{2} \klb*{\frac{-\powe{k\beta} - 1}{\i\omega_n^\F + k} - \frac{-\powe{-k\beta} - 1}{\i\omega_n^\F - k} - \powe{k\beta} \frac{-\powe{-k\beta} - 1}{\i\omega_n^\F - k} + \powe{-k\beta} \frac{-\powe{k\beta} - 1}{\i\omega_n^\F + k}}\\
&= \frac{\delta_{\sigma,-\sigma'}}{4k\cosh^2\kla*{\frac{k\beta}{2}}} \frac{1}{2} \klb*{2k \frac{1 + \powe{k\beta}}{k^2 - \kla*{\i\omega_n^\F}^2} + 2k \frac{1 + \powe{-k\beta}}{k^2 - \kla*{\i\omega_n^\F}^2}}\\
&= \frac{\delta_{\sigma,-\sigma'}}{k^2 - \kla*{\i\omega_n^\F}^2}
\end{align}
which is the expected result.
Note that it is crucial to respect the time-ordering.\\
In real time we would expect the calculation to be somewhat easier, but instead we encounter a conceptual problem right away.
If we continue in analogy to the Matsubara calculation we will end up with zero, since the anti-commutator will not yield any off-diagonal term here.
Instead, let us consider the causal Green's function for $t>0$ \TODO{TODO!}
\begin{align}
G_{1,\sigma\sigma'}^\F(k,t) &= (-\i)^2 \sum_{qs} \intr{ \mathcal{T} \bra*{c_{k\sigma}(t) c_{k,-\sigma}^\dagger(0) c_{qs}^\dagger(t') c_{q,-s}(t')}_0}{t'}\\
&= -\delta_{\sigma,-\sigma'} \powe{-\i\sigma kt} \bra*{c_{k\sigma} c_{k,-\sigma}^\dagger c_{k\sigma}^\dagger c_{k,-\sigma}}_0 \intx{0}{t}{ \powe{2\i\sigma kt'} }{t'} + \dots\\
&= \delta_{\sigma,-\sigma'} \powe{-\i \sigma kt} \frac{-1}{1 + \powe{-\sigma k\beta}} \frac{1}{\powe{-\sigma k\beta} + 1} \frac{\powe{2\i\sigma kt} - 1}{2\i\sigma k} + \dots\\
&= \frac{\delta_{\sigma,-\sigma'}}{\kla*{1 + \powe{-\sigma k\beta}}^2} \frac{\sin(kt)}{k} + \dots
\end{align}
\QnA{Note that in full generality $G^\r(x,t) = \Theta(t) \klb*{G^\F(x,t) - \kla*{G^\F(-x,-t)}^\dagger}$ allows us to obtain the retarded Green function from the Feynman/causal/time-ordered version.}
\subsection{Bosonized language}
\subsubsection{Exact solution}
Within bosonization the simple interaction term becomes a highly non-trivial sine-Gordon potential.
Since we are dealing with an operator identity, an exact solution must still be possible along the lines of the fermionic calculation (\TODO{TODO: this looks nightmarish, but it \emph{must} be possible, right?})
\subsubsection{Perturbation theory}
TODO.






%\begin{thebibliography}{9}
%\bibitem{example} description, \url{https://www.exmapleurl}, day.\,month~2021.
%\end{thebibliography}

\end{document}